\section{MDL, structure selection, and streaming}
\label{sec:mdl-streaming}

\subsection{MDL through an everyday picture}

Imagine sending a data file to a colleague over a very slow channel. You can send the raw file, or a
``formula describing the file'' plus the part that does not match the formula. If the formula is too
simple, the leftover corrections are huge. If it is too complex, the formula itself is long.
\textbf{Minimum Description Length} (MDL) chooses the description with the smallest \emph{total
number of bits}:
\begin{equation}
  \bigl\langle \mathcal{X};\mathcal{F}\bigr\rangle
  =\bigl\langle \mathcal{F}\bigr\rangle
   +\bigl\langle \mathcal{X}\mid\mathcal{F}\bigr\rangle.
  \label{eq:mdl-total}
\end{equation}
This is Eq.~(10) in the paper. The first term is the model description cost; the second is the cost of
encoding the data given the model. A good model is not a perfect fit, but one that explains much data
with few bits.

\begin{keypoint}
MDL balances \textbf{fit} and \textbf{complexity}. It is used to choose latent dimensions, decide
which outliers to keep, and decide whether to add a new regime. This is why the paper calls the method
``automatic''.
\end{keypoint}

\subsection{Model description cost}

For each model $\Theta$, the paper charges for non-zero parameters. Let $|\cdot|$ denote the number of
non-zero entries, $c_F=32$ the bits per float, and $\log^*$ the universal code length for an integer.
The formulas follow one pattern:
\begin{align}
  \bigl\langle\mathcal{F}\bigr\rangle
  &=\sum_{\Theta\in\mathcal{F}}\bigl\langle\Theta\bigr\rangle,
  &
  \bigl\langle\Theta\bigr\rangle
  &=\bigl\langle\theta_d\bigr\rangle
   +\bigl\langle\theta_s\bigr\rangle
   +\bigl\langle\widehat{\mathcal{X}}_o\bigr\rangle,
  \label{eq:model-cost}\\
  \bigl\langle\mathbf{W}^{(\mathrm{key})}\bigr\rangle
  &=|\mathbf{W}^{(\mathrm{key})}|
    \bigl(\log d_k+\log k+c_F\bigr)
    +\log^*|\mathbf{W}^{(\mathrm{key})}|,
  \label{eq:cost-wkey}\\
  \bigl\langle\mathbf{W}^{(\mathrm{loc})}\bigr\rangle
  &=|\mathbf{W}^{(\mathrm{loc})}|
    \bigl(\log d_l+\log l+c_F\bigr)
    +\log^*|\mathbf{W}^{(\mathrm{loc})}|,\\
  \bigl\langle\theta\bigr\rangle
  &=\bigl\langle\mathbf{A}\bigr\rangle
   +\bigl\langle\mathcal{D}\bigr\rangle,\\
  \bigl\langle\mathbf{A}\bigr\rangle
  &=|\mathbf{A}|\bigl(\log d_k+\log d_l+c_F\bigr)
    +\log^*|\mathbf{A}|,\\
  \bigl\langle\mathcal{D}\bigr\rangle
  &=|\mathcal{D}|\bigl(\log d_k+2\log d_l+c_F\bigr)
    +\log^*|\mathcal{D}|.
  \label{eq:cost-theta}
\end{align}
Seasonal factors and outliers have analogous costs:
\begin{align}
  \bigl\langle\mathbf{S}^{(\mathrm{time})}\bigr\rangle
  &=|\mathbf{S}^{(\mathrm{time})}|
    \bigl(\log d_s+\log L_c+c_F\bigr)
    +\log^*|\mathbf{S}^{(\mathrm{time})}|,\\
  \bigl\langle\mathbf{S}^{(\mathrm{key})}\bigr\rangle
  &=|\mathbf{S}^{(\mathrm{key})}|
    \bigl(\log d_s+\log k+c_F\bigr)
    +\log^*|\mathbf{S}^{(\mathrm{key})}|,\\
  \bigl\langle\mathbf{S}^{(\mathrm{loc})}\bigr\rangle
  &=|\mathbf{S}^{(\mathrm{loc})}|
    \bigl(\log d_s+\log l+c_F\bigr)
    +\log^*|\mathbf{S}^{(\mathrm{loc})}|,\\
  \bigl\langle\widehat{\mathcal{X}}_o\bigr\rangle
  &=|\widehat{\mathcal{X}}_o|
    \bigl(\log n+\log k+\log l+c_F\bigr)
    +\log^*|\widehat{\mathcal{X}}_o|.
  \label{eq:cost-outlier}
\end{align}
You need not memorize each line. What matters: cost grows with the number of non-zero entries and axis
sizes, so increasing $d_k,d_l,d_s$ makes the model more expensive, and adding outliers also costs bits.
The symbol $n$ in the outlier cost is not redefined at the formula; the paper uses it as the stream
length/total size. This is an ambiguity, so this guide preserves the formula and flags it instead of
assigning a confident meaning.

\subsection{Data encoding cost}

The remaining term describes the data under the model. The paper uses Huffman coding and a Gaussian
noise model with mean $\mu$ and variance $\sigma^2$:
\begin{equation}
  \bigl\langle\mathcal{X}\mid\mathcal{F}\bigr\rangle
  =\sum_{x\in\mathcal{X}}
   -\log_2 p_{\mu,\sigma}
   \bigl(x-\widehat{x}_d-\widehat{x}_s-\widehat{x}_o\bigr).
  \label{eq:data-cost}
\end{equation}
This is the unnumbered equation immediately after Eq.~(10). Intuitively, if the residual
$x-\widehat{x}_d-\widehat{x}_s-\widehat{x}_o$ is small and consistent with expected noise, it costs few
bits; an unusually large residual costs many bits. The model must therefore fit without turning every
small fluctuation into a parameter.

\begin{plainwords}
MDL does not ask ``which model fits best?''. It asks ``which model gives the shortest description of
both the model and the data?''. An over-large model can fit perfectly but describe itself at length; an
over-small model is compact but leaves a long residual. The optimal structure is the balance point.
\end{plainwords}

\subsection{How are outliers and dimensions selected?}

For outliers, the paper tries setting each residual entry to zero and keeps the change if total cost
falls. This is principled sparsification, not a hand-picked threshold.

For dimensions, Algorithm~4 (\textsc{RankUpdate}) tries configurations that change one rank by one:
$(d_k-1,d_l,d_s)$, $(d_k+1,d_l,d_s)$, $(d_k,d_l-1,d_s)$, $(d_k,d_l+1,d_s)$,
$(d_k,d_l,d_s-1)$, $(d_k,d_l,d_s+1)$. Each candidate is re-estimated and the lowest-cost structure is
kept. In the paper's HTML/PDF listing, the candidate list repeats $(d_k,d_l+1,d_s)$ and omits
$(d_k+1,d_l,d_s)$; this is very likely a presentation error. This guide uses the complete
``change one rank'' idea and states the discrepancy.

\subsection{ModelUpdate: add, keep, or switch regime?}

Algorithm~3 compares the total cost of adding a candidate $\Theta'$ to the model set $\mathcal{F}$ with
the cost without it. If adding lowers cost, the candidate is accepted; otherwise the old model remains.
After acceptance, RankUpdate may adjust $d_k,d_l,d_s$.

\begin{center}
\resizebox{\textwidth}{!}{%
\begin{tikzpicture}[
  >=Latex,
  node distance=6mm and 6mm,
  stage/.style={draw=DTNavy, rounded corners=1.5pt, fill=DTMist,
    text width=3.2cm, minimum height=11mm, align=center, font=\sffamily\small},
  decision/.style={diamond, draw=DTClaret, fill=DTWarm, aspect=2.0,
    text width=2.2cm, align=center, font=\sffamily\small},
  arrow/.style={->, thick, DTTeal}
]
  \node[stage] (x) {new window\\$\mathcal{X}^{c}$};
  \node[stage, right=of x] (est) {ModelEstimation\\candidate $\Theta'$};
  \node[decision, right=of est] (cost) {cost\\decreases?};
  \node[stage, right=of cost] (add) {add $\Theta'$\\and RankUpdate};
  \node[stage, below=7mm of cost] (keep) {keep current model};
  \node[stage, below=7mm of add] (fc) {Forecast\\$L_f$ steps};
  \draw[arrow] (x) -- (est);
  \draw[arrow] (est) -- (cost);
  \draw[arrow] (cost) -- node[above,font=\small]{yes} (add);
  \draw[arrow] (cost) -- node[right,font=\small]{no} (keep);
  \draw[arrow] (add) -- (fc);
  \draw[arrow, dashed, DTNavy] (fc.south) -- ++(0,-0.5) -| (x.south);
\end{tikzpicture}%
}
\end{center}

\subsection{The full streaming algorithm}

Algorithm~1 initializes on the first window, sets $\mathcal{F}=\{\Theta\}$, and then for each time
$t_c$:
\begin{enumerate}
  \item take the current window $\mathcal{X}^{c}$;
  \item run ModelEstimation to obtain candidate $\Theta'$;
  \item run ModelUpdate to add/keep and update ranks;
  \item forecast $L_f$ steps using $\mathcal{F}$;
  \item advance to the next time point.
\end{enumerate}

\begin{keypoint}
Per-step cost does not depend on total stream length because step 1 always takes exactly $L_c$ slices.
The set $\mathcal{F}$ may grow when a new regime appears, and RankUpdate creates runtime spikes, but no
loop touches the entire history. This is the precise meaning of scalability in the paper.
\end{keypoint}

\subsection{Forecast and what the paper does not state}

Algorithm~1 calls
$\widehat{\mathcal{X}}_{t_c+L_f}=\mathrm{Forecast}(\mathcal{F})$, but the main text does not spell out
the Forecast formula. This guide therefore states only what is certain: forecasting uses the active
model in $\mathcal{F}$, and the trend can be extrapolated from the reaction--diffusion system. It does
\emph{not} invent a detailed recurrence for seasonal/outlier terms. For exact reproduction, read the
authors' code or implementation appendix.

\subsection{What does ``parameter-free'' mean?}

The paper says D-Tracker requires no parameter tuning. Read this in context: the user does not manually
choose $d_k,d_l,d_s$ or regime switch times; MDL selects them. Technical choices still exist, such as
the ODE solver, ALS iteration count, convergence tolerance, and the initialization search over
$d_k,d_l\in[2,4]$, $d_s\in[0,4]$ used in the experiments. ``No hyperparameters'' does not mean ``no
numerical settings whatsoever''.

\begin{commonmistake}
Do not say ``MDL always finds the true model''. MDL optimizes a coding criterion; if the noise
assumption is wrong, or reality is more complex than the model family, the MDL-optimal model can still
forecast poorly. It also does not replace causal validation of diffusion arrows.
\end{commonmistake}

\begin{selfcheck}
\begin{enumerate}
  \item If $d_k$ increases, which part of MDL rises, and which part may fall?
  \item Why are outliers not selected by a fixed threshold?
  \item What does ModelUpdate do when a candidate does not reduce total cost?
  \item Explain in one sentence why per-step runtime does not grow with stream length.
\end{enumerate}
\end{selfcheck}
